\section{Chạy lại bảng bề rộng}
\label{sec:ch06-be-rong}

\index{chỗ thắt!đo lại với attention}%
Mục~\ref{sec:ch05-bop-vector} bóp \tn{vector ngữ cảnh}{context vector} lại và
xem mô hình hỏng ở đâu. Bảng ấy là phép đo trung tâm của chương trước, nên phép
kiểm thẳng thắn nhất cho chương này là chạy lại đúng nó với cơ chế mới lắp vào.

Mọi thứ giữ nguyên: cùng corpus, cùng cách chia tập, cùng seed, cùng batch,
cùng 14 epoch, cùng learning rate, cùng \tn{đảo câu nguồn}{source reversal},
cùng giải mã tham lam, cùng cách chấm \tn{khớp đúng cả câu}{exact match}. Chỉ
đổi kiến trúc.

\begin{measured}{bảng bề rộng của chương trước, có và không có attention}
\begin{minted}{text}
d_hidden  model      params   loss     exact    short    long
16        fixed      9173     0.8733   0.0033   0.0068   0.0000
16        attention  15413    0.5933   0.0600   0.1216   0.0000
32        fixed      20133    0.5584   0.1433   0.2905   0.0000
32        attention  40805    0.0088   1.0000   1.0000   1.0000
64        fixed      54341    0.1429   0.5200   0.9797   0.0724
64        attention  128453   0.0019   0.9967   1.0000   0.9934

short = source of at most 11 tokens (148 of 300), long = the rest
\end{minted}

6000 cặp huấn luyện, 300 cặp test, một seed mỗi dòng.

Đo bằng \code{experiments/ch06\_width.py} tại tag \code{ch06}.
\end{measured}

\index{đo lường!tái lập số của chương trước}%
Trước khi đọc kết quả, hãy đọc phần không phải kết quả. Ba dòng \code{fixed}
trùng khớp từng chữ số với bảng chương~\ref{ch:encoder-decoder}, kể cả cột
\code{loss}: 0.8733 và 0.0033 ở bề rộng 16, 0.5584 và 0.1433 ở 32, 0.1429 và
0.5200 ở 64. Đó không phải trùng hợp và cũng không phải mục tiêu tự thân. Nó là
điều kiện để hai bảng đặt cạnh nhau được. Nếu ba dòng ấy lệch đi, dù chỉ ở chữ
số thứ tư, thì mọi chênh lệch trong bảng này còn có thể là chênh lệch của môi
trường chạy.

\subsection*{Chỗ nhảy nằm giữa 16 và 32}

\index{attention!ngưỡng bề rộng}%
Ở bề rộng 32, mô hình vector cố định được 0.1433 và mô hình attention được
1.0000. Không phải cao hơn: đúng. Ba trăm câu test, không sai câu nào. Cùng
corpus, cùng số bước cập nhật, cùng seed.

Cột \code{long} nói cùng chuyện gắt hơn. Ở bề rộng 32, mô hình vector cố định
được \emph{tròn không} trên câu dài, không sót câu nào, trong khi câu ngắn của
nó đã lên 0.2905. Đó là hình dạng mà mục~\ref{sec:ch05-bop-vector} đã mô tả:
vector hẹp chứa được câu ngắn trước. Mô hình attention ở cùng bề rộng ấy được
1.0000 trên cả hai cột.

\subsection*{Nhưng bề rộng vẫn quan trọng, chỉ là quan trọng theo kiểu khác}

\index{attention!chỗ attention không cứu được}%
Đọc bảng từ dưới lên thì dễ kết luận quá tay, nên tôi đọc từ trên xuống.

Ở bề rộng 16, mô hình attention được 0.0600. Cao hơn 0.0033 của mô hình vector
cố định gần hai mươi lần, và vẫn là một mô hình không dịch được. Cột
\code{long} của nó tròn 0. Tôi đo thêm hai bề rộng nữa ngoài bảng, 4 và 8, và
cả hai đều cho 0.0000.

Nghĩa là attention \emph{không} làm bề rộng hết quan trọng. Nó làm bề rộng hết
quan trọng \emph{đối với việc chứa câu nguồn}. Trạng thái decoder vẫn phải đủ
rộng để làm việc của nó, tức nhớ mình đang viết dở cái gì và sinh ra từ tiếp
theo, và dưới một ngưỡng nào đó thì nó không đủ rộng cho việc ấy dù có nhìn
thấy cả câu nguồn.

Chỗ đứt gãy giữa 16 và 32 sắc một cách khó chịu: 0.0600 rồi 1.0000, không có gì
ở giữa. Bảng này không nói được vì sao nó sắc đến thế, vì nó chỉ có hai điểm ở
hai bên. Bài tập cuối chương có một câu hỏi về chỗ đó.

\subsection*{Đếm tham số trước khi mừng}

\index{đo lường!so hai kiến trúc phải đếm tham số}%
Ở cùng \(d_h\), mô hình attention lớn hơn. Ở bề rộng 32 nó có 40805 tham số so
với 20133, tức gấp đôi. Nên dòng 32 so với dòng 32 không phải một phép so công
bằng, và nếu chương dừng ở đó thì nó đang trình bày \enquote{to hơn thì tốt
hơn} dưới cái tên attention.

Phép so sống sót được phản biện ấy nằm ở hai dòng \emph{không} cùng bề rộng.
Mô hình attention ở \(d_h = 32\) mang 40805 tham số và được 1.0000. Mô hình
vector cố định ở \(d_h = 128\), mà chương trước đo được, mang 171909 tham số và
được 0.8533 với câu dài ở 0.7105. Ít hơn bốn phẩy hai lần tham số, cao hơn ở
cột \code{exact} và cột \code{long}, và bằng nhau ở cột \code{short} vì cả hai
đều đã kịch trần 1.0000 ở đó.

Đó là phát biểu tôi thấy đứng vững: trên corpus này, một mô hình attention nhỏ
dịch đúng nhiều hơn một mô hình vector cố định lớn gấp bốn lần nó. Không phải
nhờ nó lớn, vì nó nhỏ hơn.

\subsection*{Một dòng đi xuống, và nó không phải bug}

\index{đo lường!chênh lệch nhỏ hơn độ tản seed}%
Ở bề rộng 64, mô hình attention được 0.9967 trong khi ở bề rộng 32 nó được
1.0000. Rộng hơn mà thấp hơn.

Chênh lệch ấy là \emph{một câu} trong ba trăm. Và mục~\ref{sec:ch05-giai-ma} đã
đo được rằng ba mô hình chỉ khác nhau ở seed cho 0.8733, 0.6767 và 0.8700, tức
tản gần hai mươi điểm. Một câu trên ba trăm nằm sâu dưới mức ấy. Bảng này không
phân biệt được 0.9967 với 1.0000 và không nên cố.

Cái nó phân biệt được là 0.1433 với 1.0000, và 0.0724 với 0.9934. Chọn số lần
lặp theo cỡ hiệu ứng, như chương trước đã nói: hiệu ứng ở đây là cả quãng từ
gần không tới tròn một, nên một seed mỗi dòng là đủ, đúng lý do một seed mỗi
dòng đủ cho bảng chương trước.

\subsection*{Bảng này chưa nói gì về bài báo}

\index{corpus đồ chơi!giới hạn của bảng bề rộng}%
Văn phạm sinh ra corpus này hữu hạn, \tn{từ điển}{vocabulary} đóng, không có từ
\tn{ngoài từ điển}{out-of-vocabulary}, không có nhập nhằng, câu dài nhất 23
token. Bài báo chạy WMT'14 Anh-Pháp với 348 triệu từ và câu tới 50 từ. Một mô
hình đạt 1.0000 ở đây nghĩa là nó học thuộc được một văn phạm nhỏ, không nghĩa
là nó dịch được.

Cái bảng này kiểm là \emph{lập luận}: nếu chỗ thắt là tỷ lệ giữa lượng thông
tin và bề rộng chỗ nó đi qua, thì bỏ ràng buộc bề rộng đi phải làm mô hình hết
hỏng theo độ dài câu. Lập luận ấy kiểm được ở mọi quy mô mà chỗ thắt còn có
nghĩa, và ở quy mô này nó đúng.

Ở quy mô của bài, cùng phép so cho cùng chiều: RNNsearch-50 được BLEU 26.75 so
với 17.82 của RNNencdec-50, trên cùng dữ liệu và cùng quy trình huấn luyện.
Mục~\ref{sec:ch06-con-lai-gi} đọc kỹ bảng ấy.
